\def\ChapterCode{KRI}
\chapter[\CO{[KRI]} Opfer von Kriminalität, Misshandlung und Vernachlässigung] 
{Opfer von Kriminalität, Misshandlung und Vernachlässigung}


\ChapterIntro
{%
~~~
}
{%
\begin{description}
  \item[Maintainer:] Sebastian Gabriel
  \item[Autoren:] Diverse
  \item[Reviewer:] Standard-Reviewprozess
  \item[Version:] {\DocVersion} ({\DocVersionInfo})
  \item[SHA1:] {\DocGitCommit}
\end{description}

{\TxTeilkapitelAASS}
}

\Text{\TxQuerverweise}{%
\begin{itemize}
  \item Anzeigepflicht:
  \dotfill \FullPrettyRef{topic:anzeigepflicht}
  \item Anzeigerecht: \dotfill
  \FullPrettyRef{topic:anzeigerecht}
  \item Dokumentation: \dotfill
  \FullPrettyRef{topic:dokumentation-rechtlich},
  \FullPrettyRef{topic:dokumentation-medizinisch}
\end{itemize}
}{}{}{}{}


\clearpage
\section{Handlungen gegen die sexuelle Selbstbestimmung und Integrität}

\subsection{Sexuelle Gewalt, Vergewaltigung und geschlechtliche Nötigung}

\AuthNotes{eher akut, besonderer Fokus Opferbetreuung}

\Text{Vergewaltigung}{%
\label{topic:vergewaltigung}
Eine Vergewaltigung begeht, wer eine Person mit Gewalt,
durch Entziehung der persönlichen Freiheit oder durch
Drohung mit gegenwärtiger Gefahr für Leib oder Leben
 zur Vornahme oder Duldung des Beischlafes, oder einer dem
 Beischlaf
gleichzusetzenden geschlechtlichen Handlung,
nötigt.\ZFootnote{{\S}\,201 StGB}
}{}{}{}{}

\Text{Geschlechtliche Nötigung}{%
Eine Geschlechtliche Nötigung begeht, wer -- außer in den
Fällen von Vergewaltigung -- eine Person mit Gewalt oder
durch gefährliche Drohung zur Vornahme oder Duldung einer
geschlechtlichen Handlung nötigt.\ZFootnote{{\S}\,202 StGB}
}{}{}{}{}


\Text{Reaktion des Opfers}
{Eine Vergewaltigung oder sexuelle Nötigung ist meist ein
akuter Vorfall. Dementsprechend sieht man sich einem
\enquote{unvorbereiteten}, zumeist verstörtem Opfer gegenüber. Das
folgende Verhalten des Opfers kann äußerst unterschiedlich
sein und von Weinkrämpfen bis hin zur scheinbaren
Teilnahmslosigkeit reichen. Diese Reaktionen sind das
Ergebnis von Bewältigungsstrategien der menschlichen Psyche
(z.\,B. Verdrängung, etc.).
 Aus dem anfänglichen Verhalten können keine
Rückschlüsse auf das Ausmaß des seelischen Traumas
gezogen werden!}
{}
{}
{}
{}



\Text{Reaktion des Fachpersonals}
{Die Behandlung von Opfern von sexuellen Übergriffen stellt
auch für das Personal eine große psychische Herausforderung
dar. Es ist wichtig zu betonen, dass das Fachpersonal dem
Opfer die Last des Vorfalls nicht nehmen und auch nicht
wesentlich mildern kann. Ein Gefühl der \EE{eigenen Rat- und
Hilflosigkeit} ist normal. Übertriebene Ansprüche des
Personals an sich selbst schaden in der Situation der
Betreuung des Patienten!
}
{}
{}
{}
{}

\Text{Verhalten gegenüber dem Opfer}
{Ist man aufgrund der ungewohnten Situation unschlüssig, wie
man \enquote{richtig reagieren} soll, so soll man den Fokus auf
eine \EE{sachliche} und \EE{höfliche} Betreuung richten und
versuchen, auf der \EE{Sachebene} ausführlich zu
kommunizieren. 

Für Patienten in einer Ausnahmesituation ist es wichtig,
\Speech{\enquote{sich irgendwo anhalten zu können}} bzw.
\Speech{\enquote{einen Plan zu haben}}. Dementsprechend wichtig
ist auch die Information über alle zu setzenden 
Maßnahmen und das weitere Procedere (sofern absehbar). }
{}
{}
{}
{}


% M %%%%%%%%%%%%%%%%%%%%%%%%%%%%%%%%%%%%%%%%%%%%%%%%%%% M %
% Vergewaltigung und geschlechtliche Nötigung
\ProcedurePrintDefault{MT74021C}




\subsection{Sexueller Missbrauch}

\AuthNotes{eher chronisch, besonderer Fokus Erkennung und Dokumentation}

Sexueller Missbrauch umfasst grundsätzlich die Vornahme oder
Verleitung zu geschlechtlichen Handlungen unter bestimmten
Umständen, unter welchen die betroffene Person aufgrund ihres
Alters, Geschäftsfähigkeit, bestehender Abhängigkeiten oder
anderer Gründe \E{besonderen Schutz durch das Gesetz
genießt}.
Darunter fällt: 
\begin{itemize}
  \item Sexueller Missbrauch einer wehrlosen oder psychisch beeinträchtigten
  Person\ZFootnote{%
  {\S}\,205. (1) Wer eine wehrlose Person oder eine Person,
  die wegen einer Geisteskrankheit, wegen einer geistigen
  Behinderung, wegen einer tiefgreifenden Bewusstseinsstörung
  oder wegen einer anderen schweren, einem dieser Zustände
  gleichwertigen seelischen Störung \E{unfähig ist, die
  Bedeutung des Vorgangs einzusehen} {\ldots}, {\ldots}
  dadurch missbraucht, dass er an ihr eine \E{geschlechtliche
  Handlung} vornimmt oder von ihr an sich vornehmen lässt oder
  sie zu einer geschlechtlichen Handlung mit einer anderen
  Person oder, {\ldots}, dazu verleitet, eine geschlechtliche
  Handlung an sich selbst vorzunehmen, ist mit Freiheitsstrafe
  {\ldots} zu bestrafen.
  
  (2) {\ldots}
  }
  \item Schwerer sexueller Mißbrauch von
  Unmündigen\ZFootnote{% {\S}\,206. (1) Wer mit einer
  unmündigen Person den Beischlaf oder {\ldots}
  gleichzusetzende geschlechtliche Handlung unternimmt, ist
  mit Freiheitsstrafe {\ldots} zu bestrafen.
  
  (2) {\ldots} (3) {\ldots} (4) {\ldots}
  }
  \item Sexueller Mißbrauch von Unmündigen\ZFootnote{%
  {\S}\,207. (1) Wer {\ldots} eine geschlechtliche Handlung
  an einer \E{unmündigen Person} vornimmt oder von einer
  unmündigen Person an sich vornehmen läßt, ist mit
  Freiheitsstrafe {\ldots} zu
  bestrafen.
  
  (2) {\ldots} (3) {\ldots}
  
  (4) Übersteigt das Alter des Täters das Alter der unmündigen
  Person nicht um mehr als vier Jahre {\ldots}, so ist der
  Täter nach Abs. 1 und 2 \E{nicht zu bestrafen}, es sei denn,
  die unmündige Person hätte das zwölfte Lebensjahr noch nicht
  vollendet. } \item Sexueller Missbrauch von Jugendlichen
  \ZFootnote{ {\S}\,207b. (1) Wer an einer Person, die das
  \E{16. Lebensjahr} noch nicht vollendet hat und aus
  bestimmten Gründen noch nicht reif genug ist, {\ldots} unter
  Ausnützung dieser mangelnden Reife sowie seiner
  \E{altersbedingten Überlegenheit} eine geschlechtliche
  Handlung vornimmt, {\ldots}, ist {\ldots} zu bestrafen.
  
  (2) Wer an einer Person, die das 16. Lebensjahr noch nicht
  vollendet hat, \E{unter Ausnützung einer Zwangslage} dieser
  Person eine geschlechtliche Handlung vornimmt, {\ldots}, ist
  {\ldots} zu bestrafen.
  
  (3) Wer eine Person, die das 18. Lebensjahr noch nicht
  vollendet hat, unmittelbar durch ein Entgelt dazu verleitet,
  eine geschlechtliche Handlung {\ldots} vorzunehmen {\ldots},
  ist {\ldots} zu bestrafen.
  } \item Sittliche Gefährdung von Personen unter sechzehn
  Jahren\ZFootnote{ {\S}\,208. (1) Wer eine Handlung, die
  geeignet ist, die sittliche, seelische oder gesundheitliche
  Entwicklung von Personen unter sechzehn Jahren zu gefährden,
  vor einer unmündigen Person oder einer seiner Erziehung,
  Ausbildung oder Aufsicht unterstehenden Person unter
  sechzehn Jahren vornimmt, um dadurch sich oder einen Dritten
  geschlechtlich zu erregen oder zu befriedigen, ist mit
  Freiheitsstrafe bis zu einem Jahr zu bestrafen, es sei denn,
  daß nach den Umständen des Falles eine Gefährdung der
  unmündigen oder Person unter sechzehn Jahren ausgeschlossen
  ist.
  
  (2) {\ldots}
  }
  \item Mißbrauch eines Autoritätsverhältnisses\footnote{%
  {\S}\,212. (1) Wer
  
  1. mit einer mit ihm in absteigender Linie verwandten minderjährigen Person, seinem
  minderjährigen Wahlkind, Stiefkind oder Mündel oder
  
  2. mit einer minderjährigen Person, die seiner Erziehung, Ausbildung oder Aufsicht
  untersteht, unter Ausnützung seiner Stellung gegenüber dieser Person
  
  eine geschlechtliche Handlung vornimmt {\ldots}, ist {\ldots} zu bestrafen.
  
  (2) Ebenso ist zu bestrafen, wer
  
  1. als Arzt, klinischer Psychologe, Gesundheitspsychologe, Psychotherapeut,
  \EE{Angehöriger eines Gesundheits- und Krankenpflegeberufes} oder Seelsorger
  \EE{mit einer berufsmäßig betreuten Person},
  
  2. als Angestellter einer Erziehungsanstalt oder sonst als in einer
  Erziehungsanstalt Beschäftigter mit einer in der Anstalt betreuten Person oder
  
  3. als Beamter mit einer Person, die seiner amtlichen Obhut anvertraut ist,
  unter Ausnützung seiner Stellung dieser Person gegenüber eine geschlechtliche
  Handlung vornimmt {\ldots}
  }
\end{itemize}
Der sexuelle Missbrauch geschieht meist über lange Zeit und
ist nur selten ein \enquote{akutes} Problem. Wichtig hierbei ist
es, ihn zu erkennen und eine ausführliche Dokumentation
schon im Verdachtsfall anzufertigen.

\Experimental{
\section{Gewaltdelikte}

\subsection{\Z{Häusliche Gewalt}}

\Z{\TakeNotesLarge}}


\section{Missbrauch und Misshandlungen}

\subsection{Kindesmisshandlung}
\label{topic:kindesmisshandlung}

\DefinitionInPara{Eine \T{Kindesmisshandlung} ist die Ausübung von absichtlicher 
psychischer oder physischer Gewalt gegen Kinder.}
Die Misshandlung von Kindern kommt nicht selten
vor. Sie geschieht oft im Verborgenen und kommt in \EE{allen
sozialen Schichten} vor! Misshandlungen können
unterschiedlich schwer ausfallen, regelmäßig gibt es Fälle,
in denen Kinder schwere, dauerhafte Schäden davontragen oder
auch zu Tode kommen.

\Text{Hinweiszeichen}
{
\begin{itemize}
  \item Langes Intervall zwischen  Verletzungszeitpunkt und Verständigung der
  Rettung.
  \item Anamnese und Verletzungsmuster  widersprechen einander!
  \item Mehrere Verletzungen verschiedenen Alters.
  \item Erkennbarer, verletzungsverursachender Gegenstand (Schuhabdruck,
  \ldots).
  \item Zeichen der Vernachlässigung.
  \item Verletzungen bzw. Schmerzen in der Genitalregion (Achtung: Es gibt auch andere
  Ursachen als sexueller Missbrauch!).
  \item \enquote{Untypische} Verletzungen.
  \item Verhaltensstörungen.
\end{itemize}
}{
\begin{itemize}
  \item Langes Intervall Verletzung -- Hilfeersuchen.
  \item Widerspruch Anamnese -- Verletzungsmuster.
  \item Verletzungen verschiedenen Alters.
  \item Erkennbarer, verletzungsverursachender Gegenstand .
  \item Vernachlässigung.
  \item Verletzungen / Schmerzen in Genitalregion (Achtung: auch andere
  Ursachen!).
  \item \enquote{Untypische} Verletzungen.
  \item Verhaltensstörungen.
\end{itemize}
}{}{}{}


% M %%%%%%%%%%%%%%%%%%%%%%%%%%%%%%%%%%%%%%%%%%%%%%%%%%% M %
% Kindesmisshandlung
\ProcedurePrintDefault{KT74011C}



\section{Vernachlässigung}

\Text{Rechtliche Grundlagen}{%
Eine Vernachlässigung begeht, wer einem anderen, der seiner
Fürsorge oder Obhut untersteht und der das achtzehnte
Lebensjahr noch nicht vollendet hat oder wegen
Gebrechlichkeit, Krankheit oder einer geistigen Behinderung
wehrlos ist, körperliche oder seelische Qualen zufügt; bzw.
wer seine \EE{Verpflichtung zur Fürsorge oder Obhut} einem
solchen Menschen gegenüber \EE{gröblich vernachlässigt} und
dadurch, wenn auch nur fahrlässig, dessen Gesundheit oder
dessen körperliche oder geistige Entwicklung beträchtlich
schädigt.\ZFootnote{{\S}\,92 StGB: Quälen oder
Vernachlässigen unmündiger, jüngerer oder wehrloser
Personen}

}{

}{}{}{}

\Text{Konsequenz}
{Auch hier ist das \EE{Erkennen} wesentlich. Gerade der
Rettungs- und Krankentransportdienst hat oft die
Möglichkeit, eine Vernachlässigung vor Ort zu erkennen und
entsprechende Maßnahmen in die Wege zu leiten. Eine
ausführliche Dokumentation ist im Verdachtsfall sehr
wichtig!

Nebstbei muss auch nicht jede Vernachlässigung
strafrechtlich relevant sein; sie kann z.\,B. bei
\EE{Überforderung} der Angehörigen auftreten. Auch dann ist
das rasche Reagieren wichtig, damit sowohl den Angehörigen, als
auch dem Pflegling geholfen werden kann!


}
{}
{}
{}
{}


\ChapterEnd